%Chargement des paquets
\usepackage{amsmath}
\usepackage{amsfonts}
\usepackage{amssymb}
\usepackage{enumerate}
\usepackage{mathtools}
\usepackage{bbm}
\usepackage{xparse, etoolbox}
\usepackage{enumerate}
\usepackage{enumitem}
\usepackage{mathabx}
\usepackage{babel}[french]

%environment
\usepackage{amsthm}
\usepackage{keytheorems}
\usepackage{hyperref} 

%Interface théorème
\newkeytheorem{lem}[
	name = Lemme
]

\newkeytheorem{prop}[
	name = Proposition
]

\newkeytheorem{thm}[
	name = Théorème
]

\newkeytheorem{coro}[
	name = Corollaire
]

\renewcommand*{\proofname}{Démonstration}

\theoremstyle{definition}
\newtheorem{defn}{Définition}

\theoremstyle{definition}
\newtheorem*{nota}{Notation}

\theoremstyle{definition}
\newtheorem*{limi}{Liminaire}

\theoremstyle{remark}
\newtheorem{rmq}{Remarque}

\theoremstyle{plain}
\newtheorem*{fait}{Fait}


%Ensembles de nombres
\newcommand{\N} {\mathbb{N}}
\newcommand{\Ne}{\N^\ast}
\newcommand{\Z} {\mathbb{Z}}
\newcommand{\Q} {\mathbb{Q}}
\newcommand{\R} {\mathbb{R}}
\newcommand{\Rb}{\overline{\mathbb{R}}}
\newcommand{\Rp}{\R_+}
\newcommand{\Rm}{\R_-}
\newcommand{\K} {\mathbb{K}}
\newcommand{\Ket} {\mathbb{K^{\ast}}}
\newcommand{\Cx}{\mathbb{C}}

%Ensembles, tribus
\newcommand{\A}{\mathbb{A}}
\renewcommand{\P}{\mathcal{P}}
\newcommand{\T}{\mathcal{T}}
\newcommand{\F}{\mathcal{F}}
\newcommand{\C} {\mathcal{C}}
\newcommand{\ouv}{\mathcal{O}}
\newcommand{\B}{\mathcal{B}}
\newcommand{\M}{\mathcal{M}}
\newcommand{\etag}{\mathcal{E}}
\newcommand{\etagOp}{\etag^{+}(\Omega)}
\newcommand{\etagO}{\etag(\Omega)}
%%Lp avant quotientage
\newcommand{\Lic}{\mathcal{L}^1}
\newcommand{\Lpc}{\mathcal{L}^p}
\newcommand{\Linfc}{\mathcal{L}^{\infty}}
\newcommand{\Lifull}{\Lic(\Omega, \B, m)}
\newcommand{\Lpfull}{\Lpc(\Omega, \B, m)}
\newcommand{\Linffull}{\Linfc(\Omega, \B, m)}
%%Lp après quotientage
\newcommand{\Li}{L^1}
\newcommand{\Ld}{L^2}
\newcommand{\Lp}{L^p}
\newcommand{\Linf}{L^{\infty}}
\newcommand{\Listd}{\Li(\Omega, m)} 
\newcommand{\Ldstd}{\Ld(\Omega, m)} 
\newcommand{\Lpstd}{\Lp(\Omega, m)} 

%Quelques opérateurs
\DeclareMathOperator{\codim}{codim}
\DeclareMathOperator{\rg}{rg}
\DeclareMathOperator{\tr}{tr}
\DeclareMathOperator{\vect}{Vect}
\DeclareMathOperator{\ima}{Im}
\DeclareMathOperator{\id}{Id}
\DeclareMathOperator{\sign}{sign}
\DeclareMathOperator{\ext}{ext}
\newcommand{\ind}{\mathbbm{1}}
\newcommand*{\spr}[2]{\left(\,#1\,\middle|\,#2\,\right)}
\newcommand{\interior}[1]{%
  {\kern0pt#1}^{\mathrm{o}}%
}
\DeclareMathOperator{\card}{card}
\DeclareMathOperator{\ppcm}{ppcm}

%Commandes de pure flemme
\newcommand{\veps}{\varepsilon}
\newcommand{\intab}{\int_{a}^{b}}
\newcommand{\intR}{\int_{-\infty}^{\infty}}
\newcommand{\intinf}{\int_{- \infty}^{+ \infty}}
\newcommand{\equi}{\Leftrightarrow}
\newcommand{\Kev}{$\K$-espace vectoriel }
\newcommand{\Kevs}{$\K$-espaces vectoriels }
\newcommand{\Rev}{$\R$-espace vectoriel }
\newcommand{\Revs}{$\R$-espaces vectoriels }
\newcommand{\Cev}{$\Cx$-espace vectoriel }
\newcommand{\Cevs}{$\Cx$-espaces vectoriels }
\newcommand{\evn}{(E, \norm{.})}
\newcommand{\interf}[2]{[#1,#2]}
\newcommand{\limb}{\lim\limits_}
\newcommand{\lebn}{\lambda_n}
\newcommand{\pp}{\text{~presque partout}}
\newcommand{\derivp}[2]{\frac{\partial #1}{\partial #2}}
\newcommand{\famiu}{(u_i)_{i \in I}}
\newcommand{\sev}{\text{~s.e.v.}}
\newcommand{\Mnp}{M_{n,p}(\K)}
\newcommand{\Mn}{M_{n}(\K)}
\newcommand{\tq}{\text{~t.q.~}}
\newcommand{\mq}{~montrer~que~}
\newcommand{\et}{\quad \text{et} \quad}
\newcommand{\ou}{\text{~ou~}}
\newcommand{\Kx}{\K[X]}


%Commandes de gars chelou
\renewcommand{\subset}{\subseteq}

%Commande restriction, ty duckduckgo
\def\restr#1#2{\mathchoice
              {\setbox1\hbox{${\displaystyle #1}_{\scriptstyle #2}$}
              \restraux{#1}{#2}}
              {\setbox1\hbox{${\textstyle #1}_{\scriptstyle #2}$}
              \restraux{#1}{#2}}
              {\setbox1\hbox{${\scriptstyle #1}_{\scriptscriptstyle #2}$}
              \restraux{#1}{#2}}
              {\setbox1\hbox{${\scriptscriptstyle #1}_{\scriptscriptstyle #2}$}
              \restraux{#1}{#2}}}
\def\restraux#1#2{{#1\,\smash{\vrule height .8\ht1 depth .85\dp1}}_{\,#2}} 

%Quelques normes
\DeclarePairedDelimiter\abs{\lvert}{\rvert}
\DeclarePairedDelimiter\labs{\bigl\lvert}{\bigr\rvert}
\DeclarePairedDelimiter\norm{\lVert}{\rVert}
\DeclarePairedDelimiterXPP\normi[1]{}\lVert\rVert{_1}{\ifblank{#1}{\:\cdot\:}{#1}}
\DeclarePairedDelimiterXPP\normd[1]{}\lVert\rVert{_2}{\ifblank{#1}{\:\cdot\:}{#1}}
\DeclarePairedDelimiterXPP\normp[1]{}\lVert\rVert{_p}{\ifblank{#1}{\:\cdot\:}{#1}}
\DeclarePairedDelimiterXPP\normq[1]{}\lVert\rVert{_q}{\ifblank{#1}{\:\cdot\:}{#1}}
\DeclarePairedDelimiterXPP\norminf[1]{}\lVert\rVert{_\infty}{\ifblank{#1}{\:\cdot\:}{#1}}

%Bracket en angle
\DeclarePairedDelimiter\angl{\langle}{\rangle}

%Set, parenthèses
\newcommand{\set}[1]{\{ #1 \}}
\newcommand{\bigset}[1]{\bigl\{ #1 \bigr\}}
\newcommand{\Bigset}[1]{\Bigl\{ #1 \Bigr\}}
\newcommand{\bigpar}[1]{\bigl( #1 \bigr)}
\newcommand{\Bigpar}[1]{\Bigl( #1 \Bigr)}

%Opitmisation des opérateurs de comparaison
\DeclarePairedDelimiter\tio{o (}{)}
\DeclarePairedDelimiter\gro{O (}{)}


\pagestyle{fancy}

\fancyhead[C]{Réduction des isométries}
\fancyhead[R]{}

\begin{document}

\underline{Source :} Algèbre - Gourdon - page 268 (ou voir Grifone, exercice 28 page 259) \newline

\underline{Leçons :}
\begin{itemize}
\item 101 : Groupe opérant sur un ensemble. Exemples et applications.
\item 106 : Groupe linéaire d’un espace vectoriel de dimension finie E, sous-groupes de GL(E). Applications.
\item 161 : Espaces vectoriels et espaces affines euclidiens : distances, isométries.
\end{itemize}

\begin{nota}
Par défaut, le produit scalaire (éventuellement hermitien) est dénoté par $\angl{\cdot, \cdot}$.
\end{nota}

\begin{lem}
\label{lem:stab}
Soit $E$ un espace euclidien (resp. hermitien) et $u$ une isométrie (resp. un endomorphisme unitaire). 
Si $F$ est un s.e.v. de $E$ stable par $u$, alors $F^{\perp}$ est stable par $u$.
\end{lem}

\begin{proof}
Soit $x \in F^{\perp}$ fixé, on veut montrer que, pour tout $y \in F, \angl{u(x), y} = 0$.
Comme $F$ est stable par $u$ et que $\ker{u} = \set{0}$ on déduit que $F$ et $u(F)$ sont isomorphes.
Ainsi, pour $y \in F$, on exhibe $x' \in F \tq y = u(x')$ puis on écrit :
\[ \angl{u(x), y} = \angl{u(x), u(x')} = \angl{x, x'} = 0 . \] 
Ce qui achève la démonstration. 
\end{proof}

\begin{thm}
Soit $E$ un espace euclidien et $u$ une isométrie de $E$. Alors, il existe une base orthonormale $\B$ de $E$ 
dans laquelle la matrice de $u$ a la forme par blocs :
\begin{equation}
M_{\B}(u) = \begin{pmatrix}
R(\theta_1) & & & & & \\
 & \ddots & & & 0 & \\
 & & R(\theta_r) & & & \\
 & & & \veps_1 & & \\
 & 0 & & & \ddots & \\
 & & & & & \veps_s
\end{pmatrix}
\end{equation}
où, pour tout $j, \veps{j} \in \set{-1,1}$ et, pour tout i,
\[ R(\theta_i) = \begin{pmatrix} \cos{\theta_i} & - \sin{\theta_i} \\ \sin{\theta_i} & \cos{\theta_i} \end{pmatrix} \in M_2(\R), \quad
\text{avec} \quad \theta_i \in \R, \theta_i \not \equiv 0 \pmod \pi. \]
\end{thm}

\begin{proof}
On procède par récurrence sur $n = \dim{E}$. Pour $n=1$, c'est évident. \\
Supposons le résultat vrai jusqu'au rang $n-1 \geq 1$ quelconque et démontrons qu'il reste vrai au rang $n$. 
On sépare la démonstration en deux cas.
\begin{enumerate}[label=(\alph*)]

\item
Si l'isométrie $u$ admet au moins une valeur propre $\lambda$ réelle, alors en exhibant $x$ un vecteur propre associée à $\lambda$
on a :
\[ \norm{u(x)} = \norm{x} = \norm{\lambda x} = \abs{\lambda} \norm{x} , \]
et donc $\lambda \in \set{-1, 1}$.
De plus, $F = \vect{x}$ est stable par $u$ donc, d'après le lemme $\ref{lem:stab}$, il en est de même pour $F^{\perp}$.
Ainsi, $\restr{u}{F^{\perp}}$ est une isométrie de $F^{\perp}$, on peut donc lui appliquer l'hypothèse de récurrence et exhiber une base
orthonormale $\B_0$ de $F^{\perp}$ dans laquelle la matrice de $\restr{u}{F^{\perp}}$ a la forme (1).
La base $\B = \B_0 \cup \set{x/\norm{x}}$ est orthonormale et, dans cette base, la matrice de $u$ a aussi la forme attendue.

\item 
Si l'isométrie $u$ n'a aucune valeur propre réelle, on définit l'endomorphisme $v$ par $v=u+u^{\ast}$.
Comme $u$ est autoadjoint, il est diagonalisable, en particulier toutes ses valeurs propres sont réelles c'est le théorème spectral).
Notons $\lambda$ une telle valeur propre et $x$ un vecteur propre associé, on a donc $u(x) + u^{\ast}(x) = \lambda x$, soit :
\begin{align*}
u(u(x) + u^{\ast}(x)) & = u^2(x) + x \\
& = \lambda u(x) ,
\end{align*}
autrement dit :
\begin{equation}
u^2(x) = \lambda u(x) -  x .
\end{equation}
Par ailleurs, la famille $(x, u(x))$ est libre puisque $u$ n'admet pas de valeur propre réelle. 
On pose $F = \vect(x, u(x))$, $F$ est de dimension 2 et est stable par $u$ selon (2).
Notons $A = \begin{pmatrix} a & b \\ c & d \end{pmatrix}$ la matrice de l'isométrie $\restr{u}{F}$ dans une base orthonormale $\B_0$
de F. Comme $A A^{T} = A^{T} A = I_n$ on déduit que :
\[ \left \{ \begin{array}{c}
a^2 + b^2 = 1 \\
ac + bd = 0 \\
c^2 + d^2 = 1 \\
a^2 + c^2 = 1 \\
ab + cd = 0 \\
b^2 + d^2 = 1 
\end{array} \right . \]
En combinant la première et la quatrième de ces équations, on a $b^2 = c^2$ donc $b = \pm c$.
Or $b=c$ est impossible, sinon $A$ serait symétrique et les valeurs propres de $\restr{u}{F}$ seraient réelles 
(toujours le théorème spectral !). De même, $b=c=0$ est impossible, donc $b=-c \neq 0$.
En remplaçant, par exemple dans la deuxième équation, on trouve $a=d$.
Enfin, comme $a^2 + c^2 = 1$, il existe un réel $\theta \in \R$ tel que :
\[ \left \{ \begin{array}{c c}
a & = \cos{\theta} \\
c & = \sin{\theta}
\end{array} \right . , \]
autrement dit, la matrice $A$ est de la forme :
\[ A = \begin{pmatrix} \cos{\theta} & - \sin{\theta} \\ \sin{\theta} & \cos{\theta} \end{pmatrix} , \]
avec de plus $\theta \not \equiv 0 \pmod \pi$ car $b \neq 0$. \\
On finit la démonstration en notant que, d'après le lemme $\ref{lem:stab}$, $F^{\perp}$ est stable par $u$, on peut donc appliquer
l'hypothèse de récurrence à l'endomorphisme $\restr{u}{F^{\perp}}$ et exhiber une base orthonormale $\B_1$ dans laquelle
la matrice de $\restr{u}{F^{\perp}}$ a la forme (1). On conclut en concatenant $\B_0$ et $\B_1$.
\end{enumerate}
\end{proof}

On peut aussi énoncer une version matricielle de ce théorème.

\begin{thm}
Soit $M \in M_n(\R)$ une matrice orthogonale. Alors il existe une matrice orthogonale $P$ telle que $P^{T}AP$ ait la forme (1).
\end{thm}

\end{document}